\documentclass[10pt,a4paper]{article}
\usepackage[a4paper,margin=1cm,top=0.6cm,bottom=0.6cm]{geometry}
\usepackage{fontspec}
\usepackage[brazilian]{babel}
\usepackage{xcolor}
\usepackage{tikz}
\usetikzlibrary{calc, positioning, arrows.meta}
\usepackage{amsmath, amssymb}
\usepackage{booktabs}
\usepackage{multicol}
\usepackage{microtype}

% --- TIPOGRAFIA (Estilo Crisis) ---
\setmainfont{Helvetica Neue}[
  BoldFont={Helvetica Neue Bold},
  ItalicFont={Helvetica Neue Italic},
  Scale=0.95
]
\newfontfamily\headerfont{Helvetica Neue Condensed Bold}[Scale=1.4]
\newfontfamily\idfont{Helvetica Neue Condensed Bold}[Scale=2.2]
\newfontfamily\idlabelfont{Helvetica Neue Bold}[Scale=0.6]

% --- CORES ---
\definecolor{crisisBlue}{RGB}{0, 150, 210} 
\definecolor{crisisOrange}{RGB}{230, 130, 40}
\definecolor{crisisGray}{RGB}{60, 60, 60}
\definecolor{crisisLightBlue}{RGB}{235, 245, 252}
\definecolor{crisisRed}{RGB}{180, 40, 40}
\definecolor{crisisLightRed}{RGB}{252, 235, 235}

\pagestyle{empty}
\setlength{\parindent}{0pt}

% --- MACROS DE LAYOUT ---

\newcommand{\crisisHeader}[4]{% #1=cor, #2=título, #3=subtítulo, #4=ID
    \begin{tikzpicture}[remember picture, overlay]
        \fill[#1] ($(current page.north west)+(0,0)$) rectangle ($(current page.north east)+(0,-3.2cm)$);
        
        % Título
        \node[anchor=west, white] at ($(current page.north west)+(0.8cm,-1.6cm)$) {
            \begin{minipage}{12cm}
                {\fontsize{34}{38}\headerfont\MakeUppercase{#2}}\\[2pt]
                {\fontsize{18}{22}\headerfont\MakeUppercase{#3}}
            \end{minipage}
        };

        % Círculo ícone
        \draw[white, line width=1.5pt] ($(current page.north east)+(-5.5cm,-1.6cm)$) circle (0.8cm);
        \node[white] at ($(current page.north east)+(-5.5cm,-1.6cm)$) {
            \begin{tikzpicture}[scale=0.5, white]
                \draw[line width=1.5pt] (-0.3,0.3) arc (120:240:0.6) -- (0.3,-0.3) arc (-60:60:0.6) -- cycle;
                \draw[line width=1pt] (-0.1,0) -- (0.1,0); 
            \end{tikzpicture}
        };

        % ID BOX
        \fill[white, rounded corners=3pt] ($(current page.north east)+(-3.5cm,-0.4cm)$) rectangle ($(current page.north east)+(-0.3cm,-2.8cm)$);
        \node[#1, anchor=center] at ($(current page.north east)+(-1.9cm,-1.0cm)$) {\idfont #4};
        \node[#1, anchor=center] at ($(current page.north east)+(-1.9cm,-1.8cm)$) {\idlabelfont CRISIS PROTOCOLS};
        \fill[#1] ($(current page.north east)+(-3.5cm,-2.1cm)$) rectangle ($(current page.north east)+(-1.95cm,-2.7cm)$);
        \fill[#1] ($(current page.north east)+(-1.85cm,-2.1cm)$) rectangle ($(current page.north east)+(-0.3cm,-2.7cm)$);
        \node[white, anchor=center] at ($(current page.north east)+(-2.72cm,-2.4cm)$) {\idlabelfont SECTION ONE};
        \node[white, anchor=center] at ($(current page.north east)+(-1.07cm,-2.4cm)$) {\idlabelfont SECTION TWO};
    \end{tikzpicture}
    \vspace{3.0cm}
}

% Box de alerta estilo Crisis
\newcommand{\crisisAlert}[2]{% #1=título, #2=corpo
    \vspace{0.2cm}
    \begin{tikzpicture}
        \node[rectangle, draw=crisisRed, fill=crisisLightRed, rounded corners=2pt, inner sep=6pt, text width=\dimexpr\linewidth-16pt] {
            {\bfseries\small\textcolor{crisisRed}{#1}}\\[2pt]
            {\footnotesize #2}
        };
    \end{tikzpicture}
    \par\vspace{0.2cm}
}

% Box de sumário estilo Crisis
\newcommand{\crisisSummary}[2]{% #1=título, #2=corpo
    \vspace{0.2cm}
    \begin{tikzpicture}
        \node[rectangle, draw=crisisBlue, fill=crisisLightBlue, rounded corners=2pt, inner sep=6pt, text width=\dimexpr\linewidth-16pt] {
            {\bfseries\small\textcolor{crisisBlue}{#1}}\\[2pt]
            {\footnotesize #2}
        };
    \end{tikzpicture}
    \par\vspace{0.2cm}
}

\begin{document}

\crisisHeader{crisisBlue}{Oxímetro}{de Pulso SpO\textsubscript{2}}{42\textsubscript{A}}

\textbf{AUTOR:} Nick Mark MD | \textbf{TRADUÇÃO:} Equipe ICU\par
\vspace{0.3cm}

\textcolor{crisisBlue}{\headerfont\large PRINCÍPIO E FISIOLOGIA}\par
\vspace{0.1cm}
{\small O oxímetro mede a saturação de oxigênio arterial ($SpO_2$) de forma contínua.}\par
\vspace{0.1cm}
{\small\bfseries \textcolor{crisisBlue}{1} Absorção Diferencial}: Utiliza luz \textcolor{red}{\textbf{Vermelha}} (660nm) e \textcolor{orange}{\textbf{Infravermelha}} (940nm).
\begin{itemize}
    \setlength{\itemsep}{0pt}
    \item {\small Hemoglobina Oxigenada ($HbO_2$) absorve mais Infravermelho.}
    \item {\small Hemoglobina Reduzida ($Hb$) absorve mais Vermelho.}
\end{itemize}
{\small\bfseries \textcolor{crisisBlue}{2} Fluxo Pulsátil (AC)}: Identifica o sangue arterial pelo pulso, filtrando o tecido/veias.

\vspace{0.2cm}
\begin{center}
\begin{tikzpicture}[scale=1.1, font=\sffamily\tiny]
    \fill[crisisLightBlue] (-0.5,-0.5) rectangle (6.5,3.2);
    \draw[->, crisisGray] (0,0) -- (6,0) node[right] {$\lambda$ (nm)};
    \draw[->, crisisGray] (0,0) -- (0,3) node[above] {Absorp.};
    \draw[thick, xcolor=red!80] (0.5, 2.5) .. controls (2, 2) and (4, 0.5) .. (5.5, 0.4) node[right] {$HbO_2$};
    \draw[thick, xcolor=blue!80] (0.5, 0.4) .. controls (2, 1) and (4, 2.5) .. (5.5, 2.5) node[right] {$Hb$};
    \draw[dashed, crisisGray!50] (1.1,0) -- (1.1,3); \node[above, scale=0.8] at (1.1,3) {660};
    \draw[dashed, crisisGray!50] (4.8,0) -- (4.8,3); \node[above, scale=0.8] at (4.8,3) {940};
\end{tikzpicture}
\end{center}

\crisisAlert{ALERTA DE SEGURANÇA: HIPOXEMIA OCULTA}{
    O oxímetro pode \textbf{superestimar} a saturação em pacientes com pele escura.
    \begin{itemize}
        \setlength{\itemsep}{0pt}
        \item Risco 3x maior de hipoxemia não detectada.
        \item SEMPRE correlacione com a clínica e gasometria ($SaO_2$).
    \end{itemize}
}

\textcolor{crisisBlue}{\headerfont\large HEMOGLOBINAS ANORMAIS}\par
\vspace{0.1cm}
\begin{itemize}
    \setlength{\itemsep}{0pt}
    \item {\small \textbf{Carboxihemoglobina ($COHb$)}: Falsamente interpretada como $HbO_2$. Valores normais mesmo em intoxicação severa por CO (incêndios).}
    \item {\small \textbf{Metemoglobina ($MetHb$)}: Leitura "travada" ao redor de \textbf{85\%}, ignorando o valor real.}
\end{itemize}

\vspace{0.3cm}
\textcolor{crisisBlue}{\headerfont\large ANÁLISE DA ONDA PLETISMOGRÁFICA}\par
\vspace{0.1cm}
{\small Valide o sinal observando a onda antes de agir sobre o valor númerico.}\par
\begin{itemize}
    \setlength{\itemsep}{0pt}
    \item {\small \textbf{Sístole}: Pico da onda. \textbf{Incisura Dicrótica}: Fechamento da valva aórtica.}
\end{itemize}
\vspace{0.1cm}
\begin{center}
\begin{tikzpicture}[scale=1.2, font=\sffamily\tiny]
    \draw[gray, opacity=0.2] (0,0) grid (5,1.2);
    \draw[thick, crisisBlue] (0,0.2) .. controls (0.2,0.2) and (0.3,1.1) .. (0.5,1.1) 
         .. controls (0.7,1.1) and (0.8,0.5) .. (0.9,0.6) 
         .. controls (1.0,0.7) and (1.4,0.2) .. (1.8,0.2)
         .. controls (2.0,0.2) and (2.1,1.1) .. (2.3,1.1)
         .. controls (2.5,1.1) and (2.6,0.5) .. (2.7,0.6) 
         .. controls (2.8,0.7) and (3.2,0.2) .. (3.6,0.2);
    \node[crisisBlue, font=\bfseries] at (0.5,1.3) {SÍSTOLE};
    \node[crisisGray] at (1.2,0.7) {Dicrótica};
\end{tikzpicture}
\end{center}

\crisisSummary{SUMÁRIO TÉCNICO \& REGRA DE OURO}{
    \begin{tabular}{@{}lp{10cm}@{}}
        \textbf{Choque} & Reduz sinal $\rightarrow$ Erro. Escolha probe central (testa/orelha). \\
        \textbf{Interferência} & Esmalte escuro, luz intensa, movimento excessivo. \\
        \textbf{Lag Time} & Atraso de 15--60s na detecção periférica. \\
        \textbf{REGRA DE OURO} & Na dúvida, o padrão-ouro é a \textbf{Gasometria Arterial}. \\
    \end{tabular}
}

\vfill
\begin{center}
    {\fontsize{7}{9}\selectfont \textcolor{crisisGray}{© 2025 ICU Crisis Protocols - Adaptado de Nick Mark MD. Uso educacional.}}
\end{center}

\end{document}
